\section{The two conductor-twenty-four quartic companion class sectors}
\label{sec:quadratic8-trivial-quartic-companion-classes}

\paragraph{Scope and conventions.}
Put $p=269$, $K=\mathbb Q(\omega)$, and let
$\psi=\omega_p^{67}$ be the primitive odd quartic character
with $\psi(2)=-i$ and $i\equiv82\pmod p$.
Write $\delta_s=\chi_{8s}$ for $s=1,-1$ and
$\eta_K=\chi_{-3}$.
The fields $F_s$ of $\delta_s\psi$ are exactly those of
AC-G158-class.  In particular their $p$-primary class groups vanish.
This result supplies the missing companion class inputs for the
cubic-trivial quartic residue cases.  It asserts no strict
cohomology vanishing, local-unit nonvanishing, or actual-window
estimate on its own.

\begin{theorem}[AC-G164-class: the two quartic companion class inputs]
\label{AC-G164-class}
For $s=1,-1$ define
\[
 \begin{gathered}
 F'_s=\mathbb Q(z_s),\qquad
 z_s^4-6spz_s^2+900p=0,\qquad
 N_s=KF_s=KF'_s .
 \end{gathered}
 \tag{Q8TC.1}
\]
Then $F'_s$ is the cyclic quartic field of
$\eta_K\delta_s\psi$; these characters are respectively
$\chi_{-24}\psi$ and $\chi_{24}\psi$.
Both have primitive conductor $f=24p=6456$, quadratic subfield
$\mathbb Q(\sqrt p)$, and discriminant
\[
        D_{F'_s}=576p^3=11211902784.
 \tag{Q8TC.2}
\]
The field $F'_1$ is totally real.  Its three explicitly specified
units below have finite index $m$ modulo torsion in the full unit
group, and
\[
                     h_{F'_1}=2m,\qquad p\nmid m.
 \tag{Q8TC.3}
\]
The field $F'_{-1}$ is totally imaginary, and an explicit integral
Stickelberger annihilator proves its $p$-primary class group zero.
Consequently, for both signs,
\[
                      \operatorname{Cl}(F'_s)[p]=0.
 \tag{Q8TC.4}
\]
Let $b$ generate the quartic factor in
$\operatorname{Gal}(N_s/\mathbb Q)=C_4\times C_2$, acting
trivially on $K$ and with quartic character value $-i$.
Then the exact class-sector conclusion needed over $N_s$ is
\[
       \frac{1-b^2}{2}
       \bigl(\operatorname{Cl}(N_s)\otimes\mathbb Z_p\bigr)=0.
 \tag{Q8TC.5}
\]
No vanishing of the remaining class sectors of $N_s$ and no
complete fundamental-unit basis is asserted.
\end{theorem}

\begin{proof}
\proofstep{1. Exact fields, characters and automorphisms.}
Retain the exact Gaussian-period identity already checked in
AC-G147-class, AC-G156-class and AC-G158-class:
\[
 \begin{gathered}
 v^4+pv^2+25p=0,\qquad
 v^2=-\frac{p+13\sqrt p}{2},\\
 g(v)=-\frac{v^3+219v}{65},\qquad g^2(v)=-v,
                  \qquad\psi(g)=-i .
 \end{gathered}
 \tag{Q8TC.6}
\]
Thus $w_s=\sqrt{2s}\,v$ gives the actual field $F_s$ of
$\delta_s\psi$, and $z_s=\sqrt{-3}\,w_s$ gives the field
of $\eta_K\delta_s\psi$.
Its displayed polynomial is Eisenstein at $p$, since its constant
term is $900p$ and $p\nmid900$; hence it has degree four.
The common kernel description in the quadratic-times-quartic
compositum proves the stated character identification, rather
than inferring it from the polynomial alone.
The two identities
$\eta_K\delta_1=\chi_{-24}$ and
$\eta_K\delta_{-1}=\chi_{24}$ give the claimed parities.

The primitive conductors of the two quadratic characters are
$24$, coprime to the conductor $p$ of $\psi$.
Their products are therefore primitive of conductor $24p$.
For an additional exact finite check, the units
$3229,2153,49$ modulo $6456$ are congruent to $1$ modulo
$f/2,f/3,f/p$, respectively, and have value $-1$ for both
product characters.  Thus none of the three prime factors can
be removed from the modulus.
Their squares are $\tau_p=\psi^2$.
The abelian conductor--discriminant formula applied to
$1,\tau_p,\eta_K\delta_s\psi,
\overline{\eta_K\delta_s\psi}$ gives $p(24p)^2$.
This formula, as in the previous class certificates, is a
retained standard input; no PARI discriminant output is its proof.

The exact automorphism and quadratic generator are
\[
 \begin{gathered}
 b(z_s)=\frac{s z_s^3-1314z_s}{390},\qquad
 b^2(z_s)=-z_s,\\
 a_s=\frac{z_s^2-3sp}{39s},\qquad
 a_s^2=p,\qquad b(a_s)=-a_s.
 \end{gathered}
 \tag{Q8TC.7}
\]
Polynomial reduction verifies every identity in (Q8TC.7).
The cyclotomic element $5113\equiv1\pmod{24}$,
$5113\equiv2\pmod p$ supplies this same generator, so the
quartic value is $-i$ for both signs.
For $s=1$ the four real embeddings have $z_1$ values
\[
       a,-a,c,-c,\qquad
       a=\sqrt{3(p+13\sqrt p)},\quad
       c=\sqrt{3(p-13\sqrt p)}.
 \tag{Q8TC.8}
\]
Both radicands are positive since $p>169$.
For $s=-1$ all four roots are nonreal.
The unique quadratic subfield is $\mathbb Q(\sqrt p)$,
so it is different from $K$.  Thus $N_s$ has degree eight,
and adjoining $\sqrt{-3}$ to either $F_s$ or $F'_s$ gives
the same field.

\proofstep{2. Three integral units and their $p$-saturation.}
In $F'_1$, abbreviate $z=z_1$ and take
\[
 \begin{aligned}
 U_1&=(-837+5z^2)/39,\\
 U_2&=(-292939446+7479912z+1773428z^2-46473z^3)/78,\\
 U_3&=(-2569373346-198666756z+1773428z^2+137399z^3)/78.
 \end{aligned}
 \tag{Q8TC.9}
\]
The exact monic relations, with coefficients in descending order,
are
\[
 \begin{array}{c|rrrrr}
 1&1&-328&26894&328&1\\
 2&1&-58370100&5749535154002&58370100&1\\
 3&1&58370100&5749535154002&-58370100&1 .
 \end{array}
 \tag{Q8TC.10}
\]
Multiplication by $z$ on the rational basis $1,z,z^2,z^3$
is an exact $4$ by $4$ matrix with
$z^4=1614z^2-242100$.
The verifier constructs each multiplication matrix from (Q8TC.9),
computes its characteristic polynomial, compares it with
(Q8TC.10), and checks the full Horner substitution to be zero.
Hence each $U_i$ is an algebraic integer with integral inverse;
the constant terms in (Q8TC.10) are all $1$.
The raw PARI expressions are only candidates for these exact
rational checks.

At $\ell=38737=144p+1$ the defining quartic has four distinct
roots in $\mathbb F_\ell$.
The certificate checks primality by complete trial division,
every polynomial and derivative value, a primitive root modulo
$\ell$, and the complete discrete-log table in its order-$p$
subgroup.  Because the polynomial is separable modulo $\ell$,
these are actual residue maps of the maximal order; the unit
denominators $39,78$ are also invertible there.
For a generator $h$ of $\mu_p(\mathbb F_\ell)$, use the maps
\[
 E_{F'_1}\longrightarrow\mathbb F_p,\qquad
       U\longmapsto\log_h(U^{(\ell-1)/p}).
 \tag{Q8TC.11}
\]
Three of the four rows on $(U_1,U_2,U_3)$ are
\[
 \begin{pmatrix}
 12&147&118\\257&151&147\\257&118&122
 \end{pmatrix},\qquad
       \det=-43371\equiv207\pmod p.
 \tag{Q8TC.12}
\]
Thus their classes are independent in $E_{F'_1}/E_{F'_1}^p$.
This field is totally real of degree four, so its torsion units
are $\{\pm1\}$ and its free unit rank is three.
Consequently these units generate a finite-index subgroup modulo
torsion and $p$ does not divide its index $m$.
This argument proves saturation only at $p$; it does not identify
the full integral unit lattice.

\proofstep{3. Finite Fourier expressions for the real class formula.}
Put $\xi=\chi_{-24}\psi$, a primitive even character of
conductor $f=6456$.
For any primitive even nontrivial Dirichlet character $\nu$ of
conductor $q$, the exact finite formula is
\[
 \left|L(1,\nu)\right|^2
   =\frac1q\left|
       \sum_{r=1}^{q-1}\overline{\nu(r)}
                     \log(2\sin(\pi r/q))\right|^2.
 \tag{Q8TC.13}
\]
Here terms with $\nu(r)=0$ vanish.
For completeness, for $0<\rho<1$ expand
$\log(1-\rho\zeta_q^a)$ by its absolutely convergent series.
The primitive finite Fourier identity gives
\[
 \sum_a\overline{\nu(a)}\log(1-\rho\zeta_q^a)
    =-\tau(\overline\nu)
                         \sum_{n\geq1}\nu(n)\rho^n/n.
 \tag{Q8TC.14}
\]
At a nonunit frequency the finite Fourier sum is zero: a unit
in the kernel of reduction to the smaller modulus, on which the
primitive character is nontrivial, leaves the exponential fixed.
At unit frequencies reindexing gives the stated Gauss sum.
Finite Fourier orthogonality then gives
$|\tau(\overline\nu)|^2=q$.
Letting $\rho\uparrow1$ is legitimate by Abel's theorem and
the bounded partial sums of a nontrivial Dirichlet character.
Pairing $a$ and $q-a$ removes the imaginary logarithm for even
$\nu$, proving (Q8TC.13).
This finite expression has no infinite-series truncation error.

For the primitive real even character $\tau_p$ the same identity
before taking absolute values reads
\[
 L(1,\tau_p)=-\frac1{\sqrt p}
       \sum_{r=1}^{p-1}\tau_p(r)\log(2\sin(\pi r/p)).
 \tag{Q8TC.15}
\]
The sign of its real Gauss sum is fixed by a rigorous positive
interval for $\sum_r\tau_p(r)\cos(2\pi r/p)$.
Its square is $p$ by the finite Gauss identity, so positivity
selects $+\sqrt p$; a square-root sign is not assumed from a
floating-point approximation.

The zeta factorization and the analytic class-number formula give
\[
 \begin{gathered}
 \mathop{\rm Res}_{s=1}\zeta_{F'_1}(s)
                  =L(1,\tau_p)|L(1,\xi)|^2,\\
 h_{F'_1}R_{F'_1}
        =\frac{\sqrt{D_{F'_1}}}{8}
                              L(1,\tau_p)|L(1,\xi)|^2.
 \end{gathered}
 \tag{Q8TC.16}
\]
The factor $8$ is $2^{r_1}(2\pi)^{r_2}/w_F$ for
$(r_1,r_2)=(4,0)$ and $w_F=2$.
The analytic class-number formula and the projected weighted-log
regulator convention are retained primary inputs, already audited
in the previous certificates.

\proofstep{4. Original unit index and exact integer isolation.}
Form the $4$ by $3$ matrix with entries
$\log|\sigma(U_i)|$ for the four embeddings (Q8TC.8), and delete
its last row.  Its absolute determinant is the regulator $R_U$
of the specified units and satisfies $R_U=mR_{F'_1}$.
All embeddings are real, so every logarithmic weight is one.
No complex-place factor two is inserted here.
The finite ball computation, at precision $512$ bits, encloses
every square root, sine, logarithm, matrix determinant and finite
sum in (Q8TC.13)--(Q8TC.16).
It verifies that all unit values and sines are nonzero, that
$R_U>0$, and that the quadratic Gauss sum is positive.
It yields, with $r=h_{F'_1}/m$,
\[
 \begin{gathered}
 R_U=4609.859649079388599074431868116469568198910\ldots,\\
 r>1,\qquad |r-2|<3\cdot10^{-139}.
 \end{gathered}
 \tag{Q8TC.17}
\]
The exact output balls, rather than the displayed decimal
abbreviation, are the evidence used in the inequalities.

Here is the required denominator bound, without any lower bound
for an unknown regulator.
Minkowski's theorem gives an integral ideal in every class with
norm at most
\[
 U=\left\lceil\frac{4!}{4^4}
                  (\lfloor\sqrt D\rfloor+1)\right\rceil=9927,
                    \qquad D=11211902784.
 \tag{Q8TC.18}
\]
If $a_n$ counts integral ideals of norm $n$, then
\[
 \begin{aligned}
 h_{F'_1}&\leq\sum_{n\leq U}a_n
       \leq U^{4/3}\zeta_{F'_1}(4/3)
       \leq U^{4/3}\zeta(4/3)^4\\
       &\leq 4^4\left\lceil\sqrt[3]{U^4}\right\rceil
         =H=54617600.
 \end{aligned}
 \tag{Q8TC.19}
\]
The Euler-factor inequality and
$\zeta(4/3)\leq1+\int_1^\infty x^{-4/3}dx=4$
justify the last estimates.
Since $r>1$, we have $m<h_{F'_1}\leq H$.
If $h_{F'_1}-2m$ were a nonzero integer, then
$|r-2|\geq1/m>1/H$.
The computed error is smaller than $1/H>1.83\cdot10^{-8}$,
so $h_{F'_1}=2m$ exactly.
Together with (Q8TC.12), this proves
$\operatorname{Cl}(F'_1)[p]=0$.
The provisional candidate integer $2$ only selected which
integer this rigorous test attempted; it was not a class-number
or completeness premise.

\proofstep{5. An integral annihilator for the odd companion.}
Put $\xi_- =\chi_{24}\psi$, so that $\xi_-$ is odd,
primitive of conductor $f=6456$, and $\xi_-(5)=1$.
The finite exact calculation is
\[
 B_{1,\xi_-}:=\frac1f\sum_{a=1}^{f-1}a\xi_-(a)
           =-38+10i,\qquad |-38+10i|^2=1544\equiv199\pmod p.
 \tag{Q8TC.20}
\]
Let $C=\mathbb Q(\zeta_f)$ and let
$\sigma_a(\zeta_f)=\zeta_f^a$.
Sharifi's Definition~4.5.1 and Lemma~4.5.2 use
\[
 \theta_C=\frac1f\sum_{(a,f)=1}a\sigma_a^{-1},\qquad
 I_C=\mathbb Z[G]\theta_C\cap\mathbb Z[G],\quad
                         G=\operatorname{Gal}(C/\mathbb Q).
 \tag{Q8TC.21}
\]
Theorem~4.5.6 states that $I_C$ annihilates $\operatorname{Cl}(C)$.
Its complete printed proof and the entire applied lemma have been
read at the source; the upstream class-field and Gauss-sum inputs
in that proof remain imported theorems.
Use the integral element
\[
 S=(5-\sigma_5)\theta_C
    =\sum_{(a,f)=1}\left\lfloor\frac{5a}{f}\right\rfloor
                                                   \sigma_a^{-1}.
 \tag{Q8TC.22}
\]
Reindexing $\sigma_5\theta_C$ gives coefficient
$\{5a/f\}$, which proves the equality in (Q8TC.22).
The source uses the negative of this generator; both signs
belong to the annihilating ideal.
On the two character lines $\xi_-$ and $\overline{\xi_-}$,
the eigenvalues of $S$ are exactly
\[
 (5-\xi_-(5))B_{1,\overline{\xi_-}}
             =-152-40i,\qquad -152+40i.
 \tag{Q8TC.23}
\]
Their reductions under $i\mapsto82$ are respectively $65,169$,
both nonzero modulo $p$.  The inverse embedding interchanges
these two units.
The finite verifier checks every floor coefficient and both
Gaussian-integer evaluations, retaining the inverse
$\sigma_a^{-1}$ convention throughout.

Now $[C:F'_{-1}]=\varphi(6456)/4=536$, prime to $p$.
Extension of ideals injects the $p$-primary class group of
$F'_{-1}$ into that of $C$, since norm after extension is
multiplication by $536$.
This injection is Galois equivariant.
The two faithful quartic character parts are therefore killed
by the units in (Q8TC.23) and vanish.
For the remaining $b^2=1$ part, norm to $\mathbb Q(\sqrt p)$
followed by extension is $1+b^2=2$.
The $p$-primary class group of this quadratic field is zero:
AC-G152-class supplies a degree-twelve field containing it with
zero $p$-class group, and extension/norm of degree six injects
its $p$-primary class group.
Thus the $b^2=1$ part vanishes as well.
Since $2$ and $4$ are units in $\mathbb Z_p$, these even and
faithful quartic parts exhaust the class group.
This proves $\operatorname{Cl}(F'_{-1})[p]=0$ without an
analytic class-number formula for the odd companion.

\proofstep{6. The precise sector over the quadratic compositum.}
Write $A=\operatorname{Cl}(N_s)\otimes\mathbb Z_p$ and let
$c_K$ be the involution fixing $F_s$.
It sends $\sqrt{-3}$ to its negative and fixes $w_s$, so
$c_Kb^2$ fixes $z_s=\sqrt{-3}w_s$.
Its fixed field is exactly $F'_s$.
By AC-G158-class and the preceding steps, norm to either fixed
field has zero $p$-primary target.
Extension after those norms therefore gives on $A$
\[
                       1+c_K=0,\qquad 1+c_Kb^2=0.
 \tag{Q8TC.24}
\]
On the summand $b^2=-1$, adding these two identities gives
$2=0$; this forces that summand to be zero.
This is exactly (Q8TC.5), and does not require a statement
about the other summand of $A$.
\end{proof}

\paragraph{Finite evidence and retained source scope.}
The companion dictionary verifies both defining polynomials,
all automorphism and quadratic-generator identities, primitive
character witnesses, the full integral floor-coefficient
annihilator, and its two exact character evaluations.
The real-field verifier checks all three rational unit matrices,
all four simple residue maps, and the rank-three minor.
The separate ball script uses only these certified units and
finite Fourier sums, with exact input hashes in its output.
Its candidate source was an un-certified \texttt{bnfinit} run;
no \texttt{bnfcertify}, complete-unit basis, or conditional class
number is used.
The character-discriminant formula, Minkowski theorem, analytic
class-number formula, and Dirichlet unit theorem retain their
previously audited source scopes.
The new primary source is Sharifi, \emph{Iwasawa Theory},
Chapter~4, Definition~4.5.1, Lemma~4.5.2 and Theorem~4.5.6,
with the exact inverse group-ring convention in (Q8TC.21).
Only the stated class inputs are concluded.
The original-coefficient strict theorem, its actual finite-label
transport, all unlisted terms of the common-window index sum,
$\mathcal R_{\rm cub}$ and H04.4 remain separate gates.
